% !TEX TS-program = latexmk
% =============================================================================
% MASTER CV - Hamouda Chekir (Français)
% Source de vérité : CV_Hamouda_Chekir_FR_1.pdf (1 page, mêmes sections).
% Adapter : copier vers hamoudachekir_cv.tex. Ne jamais écraser ce fichier.
% Dates : utiliser "à", pas de tirets cadratin.
% =============================================================================
\documentclass[10pt,a4paper]{article}

\usepackage[french]{babel}
\usepackage[utf8]{inputenc}
\usepackage[T1]{fontenc}
\usepackage{helvet}
\renewcommand{\familydefault}{\sfdefault}
\usepackage[top=0.7cm,bottom=0.65cm,left=1.1cm,right=1.1cm]{geometry}
\usepackage{xcolor}
\usepackage{titlesec}
\usepackage{enumitem}
\usepackage{hyperref}
\usepackage{microtype}

% Babel French uses an em dash as the itemize bullet. Match the English CV (•).
\AtBeginDocument{\setlist[itemize]{label=\textbullet}}
% Do not break at existing hyphens: PDF extractors drop them (scikit-learn, prohibited-object).
\exhyphenpenalty=10000
\emergencystretch=2em
\newcommand{\kw}[1]{\mbox{#1}}

\setlength{\parindent}{0pt}
\setlength{\parskip}{0pt}
\pagestyle{empty}
\raggedbottom
\linespread{0.97}\selectfont

\definecolor{primary}{HTML}{1A365D}
\definecolor{accent}{HTML}{2B6CB0}
\definecolor{muted}{HTML}{4A5568}
\definecolor{rulecolor}{HTML}{E2E8F0}

\newcommand{\TargetTitle}{Ingénieur en informatique, \kw{full-stack} et intelligence artificielle}
\newcommand{\TargetKeywords}{React · Node.js · Python · FastAPI · Microservices · Docker · MongoDB · API REST · Agile}

\hypersetup{
 colorlinks=true,
 linkcolor=accent,
 urlcolor=accent,
 pdfauthor={Hamouda Chekir},
 pdftitle={CV - Hamouda Chekir - Master FR}
}

\titleformat{\section}
 {\color{primary}\normalsize\bfseries}
 {}
 {0pt}
 {}
 [\vspace{0.3pt}{\color{accent}\titlerule[1pt]}\vspace{2pt}]
\titlespacing*{\section}{0pt}{5pt}{1pt}

\newcommand{\job}[3]{%
 \par\noindent\textbf{\color{primary}#1}\hfill{\color{muted}\small #2}\par\vspace{0.4pt}%
 {\noindent\color{accent}\textit{\small #3}}\par\vspace{1pt}%
 \begin{itemize}[label=\textbullet, leftmargin=0.85em, itemsep=0.8pt, topsep=0.4pt, parsep=0pt, partopsep=0pt]
}

\newcommand{\proj}[2]{%
 \par\noindent\textbf{\color{primary}#1}\hfill{\color{muted}\small #2}\par\vspace{0.4pt}%
 \begin{itemize}[label=\textbullet, leftmargin=0.85em, itemsep=0.8pt, topsep=0.4pt, parsep=0pt, partopsep=0pt]
}

\newcommand{\edu}[3]{%
 \par\noindent\textbf{\color{primary}#1}\hfill{\color{muted}\small #2}\par\vspace{0.3pt}%
 {\noindent\color{accent}\textit{\small #3}}\par\vspace{2pt}%
}

\newcommand{\certline}[3]{%
 \par\noindent{\small #1, {\color{accent}\textit{#2}}\hfill{\color{muted}#3}}\par\vspace{1pt}%
}

\newcommand{\contactsep}{\hspace{5pt}{\color{rulecolor}|}\hspace{5pt}}

\begin{document}

{\fontsize{16}{18}\selectfont\bfseries\color{primary}HAMOUDA CHEKIR\par\vspace{1pt}}
{\large\color{accent}\TargetTitle\par\vspace{3pt}}
{\small\color{muted}%
 hamoudachkir@yahoo.fr
 \contactsep
 +216\,55\,913\,832
 \contactsep
 Tunis, Tunisie
\par\vspace{1pt}}
{\small\color{muted}%
 \href{https://www.linkedin.com/in/hamouda-chekir-5053572b4}{linkedin.com/in/hamouda-chekir-5053572b4}
 \contactsep
 \href{https://github.com/hamoudachekir}{github.com/hamoudachekir}
\par\vspace{4pt}}

\section{PROFIL}
{\small Ingénieur en informatique diplômé de l'ESPRIT (spécialité TWIN), orienté livraison. Je conçois et mets en production des applications web fiables en m'appuyant sur le \kw{full-stack}, l'IA appliquée et les architectures microservices. Réalisation récente chez Talan : NextHire, un ATS \kw{AI-first} de \textbf{16 microservices} (React, Node.js, Python, FastAPI, MongoDB, Docker).}

\section{EXPÉRIENCE PROFESSIONNELLE}

\job{Ingénieur logiciel, projet de fin d'études}{Février 2026 à août 2026}{Talan, Tunis}
 \item Livré un ATS \kw{AI-first} (NextHire) de \textbf{16 microservices} : 1 passerelle Node.js/Express, 2 SPA React et 13 services Python FastAPI/Flask, sur architecture MongoDB partagée et frontières de services dockerisées.
 \item Automatisé le parcours d'entretien : agent vocal LLM adaptatif et module d'intégrité YOLO (vérification faciale, détection d'objets interdits) sur sessions WebRTC en direct.
 \item Implémenté un pipeline LangGraph fusionnant transcription, comportement et signaux visuels en rapports d'embauche scorés et classés ; parsing de CV automatisé (spaCy, PaddleOCR) avec contrôle \kw{anti-hallucination}.
\end{itemize}

\job{Stagiaire développeur \kw{full-stack}}{Juin 2025 à août 2025}{SmartConseil / Dafe Management, Tunis}
 \item Livré un \kw{back-office} Laravel modulaire (régions, packs, modules) en reconstruisant les workflows CRUD en Blade et JavaScript sur PostgreSQL et MongoDB.
 \item Centralisé les outils de crawling (pilotes HTML et API, grilles, modules IP) dans une codebase unique ; livré en Agile avec revues de code, tests unitaires, SonarQube et GitLab CI/CD.
\end{itemize}

\job{Stagiaire développeur web}{Juin 2024 à août 2024}{Big Deal (Club Privilèges), Tunis}
 \item Mis en service un module de codes cadeaux numériques (génération, attribution, suivi) sous Laravel et Filament, activant des campagnes promotionnelles de bout en bout.
\end{itemize}

\job{Stagiaire développeur web}{Février 2023 à mai 2023}{HOTIX, Tunis}
 \item Mis en production le site e-commerce « Consommi Tounsi \#619 » (fonctionnalités storefront en .NET, méthode Scrum, campagne de tests et déploiement).
\end{itemize}

\section{PROJETS ACADÉMIQUES}

\proj{TalentMatch, plateforme d'enchères sécurisée (IA et blockchain)}{2025}
 \item Python, FastAPI, React, \kw{Qwen2.5-VL}, Gemini 2.5, Solidity : vérification d'authenticité, estimation de prix par web scraping et smart contracts sur Ethereum et Sepolia.
\end{itemize}

\proj{Support à la décision clinique (vision par ordinateur et RAG)}{2025}
 \item Fine-tuning de modèles \kw{vision-langage} (BLIP, CLIP+LLaMA LoRA) et stack RAG (\kw{Meditron-7B}, Qdrant, PubMedBERT) pour générer des rapports médicaux sourcés à partir de radiographies.
\end{itemize}

\proj{Carnet de Santé, plateforme de suivi santé}{2025}
 \item Django, PyTorch : vision médicale ResNet50 et ResNet18 (14 pathologies thoraciques), détection d'anomalies (Isolation Forest, One-Class SVM), chatbot T5 et API sécurisée JWT.
\end{itemize}

\section{FORMATION}
\edu{Cycle ingénieur, spécialité TWIN (Technologies du Web, Internet et Réseaux)}{2023 à 2026}{ESPRIT, École Supérieure Privée d'Ingénierie et de Technologies, Tunis}
\edu{Licence en Business Computing, option E-Business}{2020 à 2023}{ESEN, École Supérieure d'Économie Numérique, Manouba}

\section{COMPÉTENCES TECHNIQUES}
{\small
\textbf{\color{primary}Langages et frameworks :} Python, JavaScript, TypeScript, Java, React, Node.js, Express, FastAPI, Flask, Django, Laravel, PHP\par\vspace{1pt}
\textbf{\color{primary}IA et machine learning :} Agents LLM, LangGraph, LangChain, RAG, PyTorch, \kw{scikit-learn}, TensorFlow, YOLO, Computer Vision, spaCy, PaddleOCR\par\vspace{1pt}
\textbf{\color{primary}Données et DevOps :} MongoDB, PostgreSQL, MySQL, Qdrant, Docker, Git, GitLab CI/CD, SonarQube, AWS, RabbitMQ, Keycloak\par\vspace{1pt}
\textbf{\color{primary}Méthodes :} API REST, architecture microservices, Scrum, Agile, tests unitaires, revue de code
}

\section{CERTIFICATIONS}
\certline{CCNA Switching, Routing and Wireless Essentials}{Cisco Networking Academy}{Septembre 2024}
\certline{AWS Academy Graduate, Cloud Foundations}{AWS Academy (badge Credly)}{Mai 2025}
\certline{Hashgraph Developer Course}{Hedera Hashgraph}{Octobre 2025}
\certline{SDG, Work Ethics and Gender Equity}{Honoris United Universities}{2024}

\section{LANGUES}
{\small Arabe : langue maternelle. Français : courant. Anglais : avancé, niveau professionnel.}

\end{document}
